\documentclass[book.tex]{subfiles}

\begin{document}

\part{Differentiable Physics}
\label{part:diffphys}

%
% Notes: SIRNN?: NN with sinusoidal act function: self-learning spectral method -- natural basis for specific kinds of applications -- lots of scales and separation in the domain
% bias and weight -- phase shift and magnitude


%In \autoref{part:physics} of this book, we reviewed physics from a programmatic, type theoretic viewpoint. 
In this part of the book, we will introduce the tools of differentiable physics. As a brief summary, differentiable physics exploits the universal approximation capability of neural networks to model aspects of physical theories.

For example, one of the first techniques we will study, Physics Informed Neural Networks (PINNs), use a neural representation to model the solution to an arbitrary partial differential equation. We shall soon find though, that the universal approximation theorem is not a panacea, since the amount of data needed to learn an arbitrary approximation may be quite high. To lower data requirements, we shall need to leveral inductive physical biases. In particular, we will leverage physical symmetries in the form of equivariances to simplify the learning task for complex physical systems. We shall also leverage ideas from computer science to construct modular intermediate "representations" of physical objects such as molecules or materials to facilitate learning of sophisticated functions.

%\textcolor{red}{TODO: Formally define inductive bias.}

\section{Making Complex Physical Systems Accessible}

One of the most powerful pedagogical aspects of differentiable physics is that it provides a tool to make previously complex physical quantities simple. Until recently, computing molecular or materials properties from the molecular formula or the unit cell was an extremely complex problem accessible only to experts. Differentiable representations have made out-of-box representations that have made these physical systems accessible to undergraduate or even high school students. 

To provide another example, physics provides many examples of very complex probability distributions that are almost infeasible to sample classically. For example, the space of chemically viable molecules, or the space of stable materials. Differentiable physics methods enable effective sampling from these and other very complex distributions. 

We orient our study of physics in this Part towards simplicity. At times, the classical physics definition is simpler so we lead with that. At others, the differentiable physics viewpoint is simpler so we start from there. Our view is that differentiable physics can aid pedagogy in physics by helping make very complex objects (such as gauge theories) accessible to simple computational analysis.


\section{Foundational Open Problems in Differentiable Physics}

From a physical perspective, work in differentiable physics is motivated by a few foundational physical problems

\begin{enumerate}
    \item \textbf{The $N$-body Problem}: Solving for the evolution of $N$ bodies under pairwise forces is an arbitrarily complex challenge. Learning a stable scaled approximation could enable solution of complex systems.
    \item \textbf{The Force Field Problem}: Force fields construct a classical approximation to complex $N$-body systems using only particle positions and atom types. Constructing a true force field that achieves quantum accuracy would enable breakthrough modeling of very complex physical systems. 
    \item \textbf{The Jastrow Factor Problem}: The Jastrow factor is an approximate term used to correct the Hartree-Fock approximation to Schrodinger's equation. Solving for the true Jastrow Factor would enable Hartree-Fock approximations to approach the accuracy of the true solution to Schrodinger's eqiuation.
    \item \textbf{The Exchange Correlation Problem}: Learning the true exchange correlation functional would enable density functional theoretic methods to be as accurate as vastly more expensive direct solutions to Schrodinger's equation.
    \item \textbf{The Fluid Closure Problem}: Solving for the true state of a fluid can involve higher and higher moments. Classically, such fluid closure techniques required approximate truncation methods. The correct fluid closure function would enable solution of arbitrarily complex fluid system.
    \item \textbf{Scaled Numerical Relativity}: The Einstein field equations yield a complex set of coupled partial differentiable quations. Numerical relativity methods allow for partial solutions of these systems in low dimensions but do not allow for scaled solution of complex gravities.
    \item \textbf{Solid Mechanics}: Writing down the stress-strain relationship for complex geometries is a challenging problem amenable to differentiable physics.
    \item \textbf{Parton, Hadron Dynamics}: Solving complex systems from quantum chromodynamics is currently beyond the capabilities of lattice quantum chromodynamic solvers.
    \item \textbf{Magnetohydrodynamic Field Solutions:} Magnetohydrodynamic systems combine the equations of the magnetic field and the fluid field. Solving complex magnetohydrodynamic geometries will prove crucial for applications in fusion
    \item \textbf{Learnable Phase Field Functionals:} Modeling interactions at the interface of two systems is often done by introducing an auxiliary ``phase field.'' Functionals on the phase field are constructed to predict meaningful properties. These functionals can potentially be modeled effectively by differentiable physics.
    %\textcolor{red}{TODO:}
\end{enumerate}
%Comments: Solid mechanics can be the next (stress-strain relationship writing down). Differential physics can address -- 
%\textcolor{red}{Governing dynamics from electromagnetism}
%- QCD (article on quanta)
%- Magnetohydrodynamics (magnetic field <-> fluid field)

Researchers have already achieved partial solutions to some of these foundational problems as we will learn in this Part, but many of these problems are still open. We hope that some of you will be inspired to solve these foundational problems in modern physics after finishing this Part.

\section{From Incomplete to Complete Physics}

One of the grand driving forces in modern physics is unification. The triumphs of theoretical physics of the last centuries have unified mechanics and optics with electromagnetism, electromagnetism with general relativity, relativity with the strong and weak nuclear forces. The remaining grand challenge is unifying general relativity with quantum mechanics.

In a parallel trend, as computers have become more capable, we have seen a parallel push towards ever more efficient algorithms to simulate larger and larger physical systems. Intriguingly, these algorithms often work through de-unification and approximation. For example, quantum effects are mostly disregarded in molecular dynamics, relativistic effects are mostly disregarded in mechanics simulators, and multibody effects are disregarded in molecular simulations. Often, these approximations are harmless and necessary, but at other times, approximations render computational tools incapable of achieving predictive accuracy due to a buildup of errors.

The dream of the differentiable physics approach is to enable us to build efficient and unified physical simulators for systems. If we could enable accurate accounting of complex, quantum, many body effects in simple simulations, we could potentially study the world through a powerful and accurate new computational lens and achieve new insights.

What problems might be amenable to differentiable physics approaches? Any time we combine two fields of physics into a joint theory we gain a more complete theory, but which likely has additional computational costs. This is potentially a place where we can use differentiable physics to achieve the accuracy of the joint theory with the efficiency of the single incomplete theories.

%\textcolor{red}{Notes: First part is "single-physics" and the second part is coupled physics. "Incomplete physics" -> "Complete physics" -- enabled by diff. physics. Field variables view of what you're solving for: v, P, mu, T, E and B -> single physics vs. coupled physics. v and T: temperature driven motion. 8*8 matrix: for combinations. Framing: Every combination leads to a specific form of incomplete physics, and the completion essentially requires a unifying piece: diff physics}



\subfile{part_diffphys_mechanics_ml}
\subfile{part_diffphys_electromagnetism_ml}
% Folded into other parts \subfile{part_diffphys_PINNs}
\subfile{part_diffphys_neural_diff_eq}
\subfile{part_diffphys_fluid_ml}
\subfile{part_diffphys_molecular_ml}
\subfile{part_diffphys_neural_potential}
\subfile{part_diffphys_deep_quantum}
\subfile{part_diffphys_materials_ml}
\subfile{part_diffphys_meshes_and_shapes}
\subfile{part_diffphys_quantum_computing_ml}
\subfile{part_diffphys_statmech_ml}
\subfile{part_diffphys_qft_ml}
\subfile{part_diffphys_hep_ml}
\subfile{part_diffphys_gr_ml}
\subfile{part_diffphys_conclusion}

\end{document}